\begin{table}[htbp]
\centering
\small
\caption{V1 E3（因果プロファイルの層状局在と共有）：残差ストリーム交換・摂動介入における因果効果ピーク深度 $d_{\text{causal}}^*$、Reader--Self間の層別重要度順位相関（Spearman $\rho_{\text{rank}}$）、および介入方向コサイン類似度 $\cos\theta$。}
\label{tab:v1_causal_profile}
\begin{tabular}{ll cccc c}
\toprule
 & & \multicolumn{2}{c}{\textbf{Peak Causal Layer}} & \multicolumn{2}{c}{\textbf{Causal Magnitude ($\times 10^{-3}$)}} & \textbf{Reader--Self Sharing} \\
\cmidrule(lr){3-4} \cmidrule(lr){5-6}
\textbf{Family} & \textbf{Variant} & Reader ($d^*$) & Self ($d^*$) & Reader (Peak) & Self (Peak) & Spearman $\rho_{\text{rank}}$ / Mean Cosine \\
\midrule
\multirow{2}{*}{Qwen 2.5 1.5B} & Base       & L16 (0.59) & L16 (0.59) & 2.56 & 2.46 & 0.805 / 0.390 \\
                       & Instruct   & L1 (0.04) & L8 (0.30) & 0.98 & 0.86 & 0.799 / 0.000 \\
\midrule
\multirow{2}{*}{Llama 3.2 1B} & Base       & L4 (0.27) & L6 (0.40) & 0.75 & 1.19 & 0.582 / 0.000 \\
                       & Instruct   & L6 (0.40) & L6 (0.40) & 2.19 & 1.92 & 0.882 / 0.414 \\
\midrule
\multirow{2}{*}{Gemma 3 1B} & Base       & L2 (0.08) & L2 (0.08) & 2.87 & 3.70 & 0.802 / 0.712 \\
                       & Instruct   & L11 (0.44) & L11 (0.44) & 4.89 & 8.21 & 0.646 / 0.266 \\
\midrule
\multirow{2}{*}{OLMo 2 1B} & Base       & L4 (0.27) & L4 (0.27) & 6.14 & 1.66 & 0.715 / 0.210 \\
                       & Instruct   & L10 (0.67) & L10 (0.67) & 0.82 & 1.17 & 0.771 / 0.000 \\
\bottomrule
\end{tabular}
\vspace{1ex}
\begin{minipage}{\linewidth}
\footnotesize
\textbf{Note:} 全モデルにおいて因果ピーク深度は Reader と Self でほぼ完全一致（同一層または近傍層）し、層別感度順位相関は $\rho_{\text{rank}} = 0.58$--$0.88$ と極めて高い。これは他者感情の認識と自己報告の生成が単に相関しているだけでなく、因果的にも共通の回路基盤に依存していることを示す。
\end{minipage}
\end{table}